\section{Simultaneous coarse estimates}\label{sec:seed}

The refinement will need an integer estimate at every interval state, including
states chosen adaptively by later computations.  We construct these estimates
as a random table, while evaluating only the entries that are requested.  The
construction builds on tree and precision-sampling ideas of Andoni,
Krauthgamer, and Onak~\cite{AndoniKrauthgamerOnak2010}; the complete bounds
needed here are proved below.

\begin{theorem}[Initial estimates]\label{thm:seed}
With the parameters and labeled interval tree of Section~\ref{sec:geometry},
there is a lazily evaluated random table $U^{(0)}(I,q)$ of integers in $[0,N]$
such that, with probability at least $0.99$, simultaneously at every state,
\begin{equation}\label{eq:seed-guarantee}
 c(I,q)\le U^{(0)}(I,q)\le A_{\mathrm{init}}c(I,q).
\end{equation}
Every zero-cost state has value zero on every execution.  Conditional on any
first request for an entry and the previously exposed randomness, its expected
evaluation cost is at most $(1+|I|)2^{o(H)}$, uniformly in the state.  This cost
bound includes arithmetic and random sampling and does not require
\eqref{eq:seed-guarantee} to hold.
\end{theorem}

\subsection{A deterministic shift tree}

The hierarchical representation below is an adaptation of the
E-distance of Andoni, Krauthgamer and Onak~\cite[Definition~3.2 and
Theorem~3.3]{AndoniKrauthgamerOnak2010}. We prove its precise bounds here,
including the treatment of unequal lengths, empty intervals and bounded
integer shifts.

Fix one state and write $n_x=|I|$, $n_y=q^+-q^-$, and $c=c(I,q)$.  If
$n_y>2n_x$, return $n_y$: the bounds $n_y-n_x\le c\le n_y$ give the required
factor two.  Otherwise
test equality directly and return zero when the strings agree.  These steps
cost $O(1+n_x)$, since in the branch requiring character comparisons
$n_y\le2n_x$.  They include the case $n_x=0$.

For the remaining case, pad the shorter string at its right end with a fresh
tagged symbol to obtain two strings of the same length $n\le2n_x$.  Denote their
distance by $d$.  The length gap is at most $c$, so $d\le2c$.  Conversely,
remove padding from any padded alignment.  A diagonal joining padding to a
real character has cost one and can be replaced by an insertion or deletion
of the same cost; other operations consuming padding can be removed.  Hence
\begin{equation}\label{eq:padding-comparison}
 1\le c\le d\le2c,\qquad 1\le n\le N,\qquad d\le n.
\end{equation}
The last bound follows by substituting characters position by position.

For each dyadic guess
\[
 k\in\{1,2,4,\ldots,2^{\ceil{\log_2 n}}\},
\]
split the first padded string in a complete balanced $B$-ary tree of depth
$D_0=\ceil{\log_B n}$.  Its leaves have width zero or one.  Allowed shifts are
the integers in $[-k,k]$.  At a width-one leaf $w$ starting at position $j$,
let $T_w(v)$ be the mismatch indicator between its character and position
$j+1+v$ of the second padded string.  A position outside the latter string
has a virtual symbol that always mismatches.  At a width-zero leaf put
$T_w(v)=0$.  At an internal node define
\begin{equation}\label{eq:coarse-tree}
 T_w(v)=\sum_{i=1}^B\min_{v'\in[-k,k]\cap\Z}
       \{T_{w_i}(v')+2|v-v'|\}.
\end{equation}
Constant shifts show that $0\le T_w(v)\le |w|$, where $|w|$ denotes the node's
width.

The virtual symbols are chosen outside the alphabets of both padded
strings. Thus every leaf value is an edit cost against a finite interval
of an extended target string, to which Lemma~\ref{lem:geometry} applies.

\begin{lemma}\label{lem:shift-tree}
At the root $o$,
\[
 d\le T_o(0),\qquad
 T_o(0)\le(1+2BD_0)d\quad\hbox{if }k\ge d.
\]
\end{lemma}
\begin{proof}
The recurrence minimizes the sum of leaf mismatches and twice the total
absolute shift difference along tree edges, with root shift zero.  For any
assignment, concatenate the shifted leaf target intervals.  By traversing
the tree in depth-first order, the sum of jumps between consecutive leaf
intervals, including the connections to the outer endpoints, is at most
twice the sum of edge shift differences.  The connection inequality from
Section~\ref{sec:geometry}, including its virtual-position version, therefore
bounds $d$ by the assignment's cost.  This proves the first assertion.

For the second, fix an optimal padded alignment and choose ordered points
$(j,b_j)$ at every source cut, with $b_0=0$ and $b_n=n$.  Assign shift $b_j-j$
to every node starting at $j$.  Its absolute value is at most $d$, because
only insertion and deletion steps change the shift.  These assignments are
allowed when $k\ge d$.

For each child, the shift difference from its parent is at most the cost of
the alignment subpath between their chosen starting points, and thus at most
the cost of the parent's subpath.  There are $B$ children per parent; the
parent subpaths at each depth partition the alignment.  Edge differences
therefore sum to at most $Bd$ at each of the $D_0$ depths.  A mismatching
width-one leaf has positive cost in its corresponding subpath: a zero-cost
subpath would be a matching diagonal.  Leaf mismatches sum to at most $d$,
and width-zero leaves contribute nothing.  The assignment has total cost at
most $(1+2BD_0)d$.
\end{proof}

\subsection{Pruning and precision sampling}

Nonuniform additive precision is the sampling principle developed by
Andoni, Krauthgamer and Onak~\cite[Lemma~3.12]{AndoniKrauthgamerOnak2010}
and refined in~\cite[Lemma~1.2]{AndoniKrauthgamerOnak2011}, with earlier
multiscale sampling antecedents in Indyk and
Woodruff~\cite{IndykWoodruff2005}. The reconstruction arguments of Andoni,
Krauthgamer and Onak already allow estimates to depend on the sampled
precisions. Here we use a finite threshold
distribution and prove an event covering all potential interval states,
together with the conditional cost of a first request.

Computing all of \eqref{eq:coarse-tree} would be too expensive. The guess
$k$ will also serve as an additive allowance at the root: on one event
of high probability, we will obtain an estimate $\widetilde T_o$ between $T_o/2-k$ and
$2T_o+k$ simultaneously for every entry and guess. Then
$2(\widetilde T_o(0)+k)$ is an upper bound on $d$ for every guess,
and the guess $d\le k<2d$ gives an $O(1+BD_0)$ approximation.
We can therefore replace a node by zero whenever its entire cost fits
within its additive allowance. Larger guesses permit more pruning;
their larger allowances are paid back in the returned estimate. Set
\[
 Q_0=H^4,\qquad \lambda=\frac1{64H},\qquad a_o=k.
\]
At a node with allowance $a_w$, independently for its children draw
$R_i$ uniformly from $\{1,\ldots,B\}$ and set
\begin{equation}\label{eq:coarse-threshold}
 t_i=\frac{a_wR_i}{BQ_0},\qquad a_{w_i}=\lambda t_i.
\end{equation}
All these random choices can be defined in advance, even below nodes that
will be pruned.  The implementation exposes only those it needs.

If $|w|\le a_w$, use the identically zero estimate $\widetilde T_w$.  At any
other leaf use its exact array.  At an unpruned internal node first compute
\[
 \widehat s_i(v)=\min_{v'\in[-k,k]\cap\Z}
       \{\widetilde T_{w_i}(v')+2|v-v'|\},
\]
and then put
\begin{equation}\label{eq:coarse-estimate}
 \widetilde T_w(v)=\sum_{i=1}^B
     \ind_{\widehat s_i(v)\ge t_i}
     \max\{\widehat s_i(v),a_w/Q_0\}.
\end{equation}
The descendants used to compute $\widehat s_i(v)$ depend on $t_i$, so
these child estimates cannot be treated as fixed in a concentration
bound. The proof below first compares their sampled contributions with
two lists formed from true child costs. It then concentrates those lists
without assuming that any child estimate is accurate.

\begin{lemma}\label{lem:coarse-precision}
Except on an event of probability less than $0.01$, simultaneously for all
initial entries, all guesses, all their nodes, and all shifts,
\begin{equation}\label{eq:coarse-accuracy}
 T_w(v)/f_w-a_w\le\widetilde T_w(v)\le f_wT_w(v)+a_w,
 \qquad
 f_w=\left(\frac{1+\lambda}{1-\lambda}\right)^{\operatorname{height}(w)}
 \le2.
\end{equation}
\end{lemma}
\begin{proof}
Fix an internal node and a shift, and write $s_i$ for the true child penalty minimum
in \eqref{eq:coarse-tree}. Let $f'$ be the factor in
\eqref{eq:coarse-accuracy} at its children's structural height. Both $s_i$
and $f'$ are independent of all thresholds. Suppose for the moment that
the child accuracy inequalities hold with this factor. Adding a nonnegative
penalty preserves the lower bound after division by $f'$ and the upper
bound after multiplication by $f'$.  Taking minima gives
\[
 s_i/f'-\lambda t_i\le\widehat s_i(v)\le f's_i+\lambda t_i.
\]
Let $u=a_w/Q_0$ and define
\[
 \phi_i(h)=\ind_{h\ge t_i}\max\{h,u\},\qquad
 h_i^- =\frac{s_i}{f'(1+\lambda)},\qquad
 h_i^+=\frac{f's_i}{1-\lambda}.
\]
If $h_i^-\ge t_i$, then
$\widehat s_i(v)\ge(1+\lambda)h_i^--\lambda t_i\ge h_i^-$.
If $\widehat s_i(v)\ge t_i$, then
$\widehat s_i(v)\le f's_i+\lambda\widehat s_i(v)$, and consequently
$\widehat s_i(v)\le h_i^+$.  These implications prove the pathwise sandwich
\begin{equation}\label{eq:threshold-sandwich}
 \phi_i(h_i^-)\le\phi_i(\widehat s_i(v))\le\phi_i(h_i^+).
\end{equation}

We now establish concentration independently of the supposition about child
accuracy.  Conditional only on ancestor thresholds, $u$ is fixed and
$t_i=uR_i/B$ are independent.  For either of the two fixed argument lists
$h_i=h_i^\pm$, if $0\le h_i<u$ then
\[
 \phi_i(h_i)=u\,\operatorname{Bernoulli}
       \left(\frac{\floor{Bh_i/u}}B\right);
\]
if $h_i\ge u$, it equals $h_i$ deterministically.  Thus
\[
 0\le h_i-\E\phi_i(h_i)\le u/B,\qquad
 \operatorname{Var}(\phi_i(h_i))\le uh_i,
 \qquad |\phi_i(h_i)-\E\phi_i(h_i)|\le u.
\]
Write $V=\sum_i h_i$ and $Z=\sum_i\phi_i(h_i)$.  The total bias is at most
$u$.  Bernstein's inequality~\cite[Theorem~4.16]{Tropp2021Probability} at
deviation $z=\lambda V+a_w/2$ gives
\[
 \Pr(|Z-\E Z|>z)
 \le2\exp\left(-\frac{z^2}{2uV+(2/3)uz}\right)
 \le2\exp(-\lambda Q_0/12).
\]
For the last inequality, if $V\ge a_w$, the denominator is at most $3uV$
and $z^2\ge\lambda a_wV$; if $V\le a_w$, the denominator is at most
$3ua_w$ and $z^2\ge a_w^2/4$.  Since $u\le a_w/2$, adding the bias yields
\begin{equation}\label{eq:threshold-concentration}
 |Z-V|\le\lambda V+a_w
\end{equation}
outside the indicated event.

These conditional bounds are uniform in the ancestor thresholds, so they
also hold unconditionally.  Take their intersection over both argument
lists and all fixed node and shift labels.  On that intersection, induction
from the leaves combines \eqref{eq:threshold-sandwich} and
\eqref{eq:threshold-concentration}.  The lower and upper bounds become
\[
 \frac{1-\lambda}{f'(1+\lambda)}T_w(v)-a_w
 \le\widetilde T_w(v)
 \le\frac{f'(1+\lambda)}{1-\lambda}T_w(v)+a_w,
\]
which is exactly the desired factor update.  A pruned node also satisfies
the inequalities because $T_w(v)\le |w|\le a_w$; exact leaves do so directly.
The factor satisfies $f_w\le\exp(2/63)<2$, since $D_0\le H$.

For completeness, the event count includes the entire domain, not merely
requested entries.  There are $O(M(N+1)^3)$ labeled main-tree states, $O(H)$
guesses per entry, $O(BN)$ nodes per coarse tree, and $O(N)$ shifts per guess.
The count of the two lists is therefore at most
\[
 CMBH(N+1)^5\le\exp(20H)
\]
for an absolute constant $C$ under the stated cutoff.  Here one may use
$B,M,H\le N$ in this regime and $\log_2(N+2)\le H$.  The union failure is at
most $2\exp(20H-H^3/768)<0.01$.  In particular, no conditioning on successful
descendants was used in the concentration calculation.
\end{proof}

\subsection{Converting the estimate to an integer seed}

For the padded entry, take the ceiling of
\begin{equation}\label{eq:initial-output}
 \min_k 2\bigl(\widetilde T_o(0)+k\bigr)
\end{equation}
and clamp it to $[0,N]$.  On the event of Lemma~\ref{lem:coarse-precision},
every candidate is at least $T_o(0)\ge d\ge c$.  A listed guess satisfies
$d\le k<2d$, and its rounded candidate is at most
\[
 4(1+2BD_0)d+4k+1
 \le (25+16BH)c
 \le BH^4c=A_{\mathrm{init}}c.
\]
Clamping preserves both bounds because $c\le N$.  Together with the direct
branches this proves the approximation assertion in Theorem~\ref{thm:seed}.

\subsection{The conditional cost of an adaptive request}

We next prove the cost assertion, including executions outside the accuracy
event.  For an array $E$ on $[-k,k]\cap\Z$, its penalty transform is computed
by a left sweep
\[
 L(v)=\min\{E(v),L(v-1)+2\}
\]
with the natural initial value, and the symmetric right sweep.  Their
pointwise minimum is
\[
 \min_{v'}\{E(v')+2|v-v'|\}.
\]
Both sweeps take $O(k)$ operations.  An identically zero child array, whose transform is also zero,
is represented by a marker without allocation.  At an unpruned internal
node the array work is $O(Bk)$, including visits and thresholds for all
children.  An unpruned leaf costs $O(k)$.  A pruned node needs only its width
test and marker.

Let $r$ be a node's depth in the coarse tree.  From
\eqref{eq:coarse-threshold}, along any fixed path,
\[
 a_w=\frac{k\prod R_e}{(64BH^5)^r},\qquad
 \E\left[\frac{a_{\mathrm{parent}}}{a_w}\,\middle|\,
          a_{\mathrm{parent}}\right]
 =64H^5\sum_{j=1}^B\frac1j\le H^8.
\]
It follows that $\E[k/a_w]\le H^{8r}$.  A node can incur array work only if
$|w|>a_w$, and pointwise
\[
 k\ind_{|w|>a_w}\le |w|\frac{k}{a_w}.
\]
Widths sum to $n$ at each depth.  Thus the expected array work for one guess,
including child overhead charged to expanded parents, is at most
\begin{equation}\label{eq:initial-array-work}
 O\bigl(1+Bn(D_0+1)H^{8D_0}\bigr).
\end{equation}
This estimate uses the random allowances alone and remains valid when the
estimated arrays are inaccurate.

All arithmetic and sampling can be performed with finite binary integers.
The quantities in \eqref{eq:coarse-threshold} are dyadic; a common denominator
for their descendants and the arrays needs
$O(D_0(p_*+5)\log H+\log H)$ bits.  Minima and maxima select existing dyadic
values, and sums preserve this common precision.  The same-shift candidate
in the penalty minimum gives
\[
 \widetilde T_w(v)\le\sum_i\widetilde T_{w_i}(v)+Ba_w/Q_0.
\]
Since $a_w\le k$ and there are $O(Bn)$ nodes, all estimate magnitudes are
bounded by $n+O(B^2nk)$.  Their bit lengths, and the cost of ordinary binary
arithmetic on them, fit the factor $H^{O(D_0+p_*+1)}$.  Each $R_i$ requires
exactly $\log_2B$ fair random bits.  Zero arrays are never materialized.

There are $O(H)$ guesses, and
\[
 D_0\le1+\frac{H}{p_*\log_2H},\qquad
 \log_2\bigl(H^{O(D_0+p_*+1)}\bigr)
 =O\left(\frac H{p_*}+(p_*+1)\log H\right)=o(H).
\]
Consequently \eqref{eq:initial-array-work}, including the arithmetic factors,
all guesses, equality tests, and padding, is $(1+|I|)2^{o(H)}$ in expectation.
Even a full coarse tree with all its shift arrays uses only polynomial
space: it has $O(Bn)$ nodes, $O(n)$ entries per array, and the bit lengths
just bounded.

Finally, assign a separate random stream to every initial entry, keyed by
its labeled main-tree node and endpoint pair, and cache the completed
answer.  Its calculation does not query other initial entries.  Before its
first request, none of the exposed computation can depend on its unused
stream: such a dependence would already have required evaluating this
entry.  Conditional on the entire preceding execution and the adaptive
request, the stream therefore has its original distribution.  The expected
cost bound applies conditionally exactly as stated; repeated requests use
the cached answer and pay only lookup work.  This deferred exposure realizes
the same full random table used in the simultaneous probability argument
and completes the proof of Theorem~\ref{thm:seed}.
